% ============================================================================
% SECCION 4: DISCUSION
% ============================================================================

\section{Discusi\'on}
\label{sec:discusion}

\subsection{AlphaFold2: Precisi\'on extraordinaria con mapeo correcto}

Este estudio revela dos perspectivas complementarias sobre la precisi\'on de AlphaFold2:

\begin{itemize}
    \item \textbf{Dataset validado} (91 prote\'inas, 5.115 pares con mapeo SIFTS expl\'icito):
    RMSD mediano de \SI{0,64}{\angstrom}, 87,7\% $< \SI{2}{\angstrom}$. Estas prote\'inas
    bien estudiadas muestran precisi\'on comparable a la variabilidad experimental.

    \item \textbf{Dataset diversificado} (10.432 prote\'inas \'unicas): RMSD mediano de
    \SI{1,26}{\angstrom}, 65,3\% $< \SI{2}{\angstrom}$. Al incluir prote\'inas m\'as
    diversas, la precisi\'on se reduce pero sigue siendo notable.
\end{itemize}

La diferencia refleja que las prote\'inas m\'as estudiadas (con muchas estructuras PDB)
tienden a ser m\'as ``f\'aciles'' para AlphaFold2, mientras que prote\'inas con menos
datos experimentales presentan mayor variabilidad.

\subsection{La importancia cr\'itica del mapeo de residuos}

Este estudio revela un problema metodol\'ogico que puede afectar a otros an\'alisis
comparativos de estructuras proteicas: \textbf{la comparaci\'on directa por n\'umero
de residuo entre archivos PDB y modelos AlphaFold produce resultados err\'oneos}.

La causa es doble:
\begin{enumerate}
    \item Los archivos PDB utilizan numeraciones de residuos que frecuentemente no
    coinciden con la secuencia can\'onica de UniProt. Esto ocurre por m\'ultiples razones:
    presencia de p\'eptidos se\~nal que se eliminan en la prote\'ina madura,
    numeraci\'on hist\'orica diferente, o convenciones de numeraci\'on espec\'ificas
    del laboratorio depositante.

    \item Las estructuras PDB multim\'ericas contienen m\'ultiples cadenas que
    corresponden a diferentes prote\'inas. Sin la informaci\'on de SIFTS sobre qu\'e
    cadena corresponde a cada identificador UniProt, se pueden comparar cadenas
    completamente diferentes.
\end{enumerate}

El efecto combinado de estos problemas transform\'o lo que deber\'ia ser un RMSD
mediano de \SI{0,64}{\angstrom}--\SI{1,19}{\angstrom} en un valor artificial de $\sim$\SI{18}{\angstrom}.
Este resultado sirve como advertencia para la comunidad: cualquier comparaci\'on
PDB--AlphaFold que no utilice el mapeo SIFTS expl\'icito de posiciones debe
considerarse con extrema precauci\'on.

\subsection{Ausencia del ``sweet spot'' de tama\~no}

An\'alisis previos (incluyendo versiones anteriores de este mismo estudio) hab\'ian
identificado un aparente ``sweet spot'' en prote\'inas de 250--300 residuos, donde
AlphaFold2 supuestamente predec\'ia con 2,05 veces mayor precisi\'on que el promedio.
\textbf{Con el mapeo SIFTS correcto, este sweet spot desaparece por completo}.

La precisi\'on de AlphaFold2 es alta y relativamente uniforme en todos los rangos
de tama\~no, con variaciones que reflejan la composici\'on del dataset (qu\'e prote\'inas
espec\'ificas est\'an representadas) m\'as que una tendencia del algoritmo. El rango
250--300 muestra RMSD medio de \SI{1,12}{\angstrom}, virtualmente id\'entico al
promedio general de \SI{1,05}{\angstrom}.

El ``sweet spot'' era un artefacto del desfase de numeraci\'on: prote\'inas cuya
numeraci\'on PDB coincid\'ia casualmente con la numeraci\'on UniProt produc\'ian RMSD
bajos, creando la ilusi\'on de un rango \'optimo de tama\~no.

\subsection{Casos con RMSD elevado}

El porcentaje de prote\'inas con RMSD $> \SI{5}{\angstrom}$ (1,4\% en el dataset validado,
16,6\% en el diversificado) corresponde a:

\begin{itemize}
    \item \textbf{Cambios conformacionales ligando-inducidos}: AlphaFold2 predice
    la conformaci\'on apo (sin ligando), mientras que muchas estructuras PDB capturan
    la prote\'ina unida a sustratos, inhibidores o cofactores. CDK2 (P24941,
    RMSD \SI{2,18}{\angstrom}) ilustra esto: las m\'ultiples conformaciones activa/inactiva
    difieren significativamente.

    \item \textbf{Prote\'inas multidominio con bisagras}: Prote\'inas con dominios
    conectados por regiones flexibles pueden adoptar diferentes orientaciones relativas,
    produciendo RMSD global elevado aunque cada dominio individual est\'e bien predicho.

    \item \textbf{Regiones intr\'insecamente desordenadas}: Segmentos que carecen
    de estructura terciaria estable pueden cristalizar en conformaciones espec\'ificas
    del empaquetamiento cristalino, diferentes de lo predicho por AlphaFold2.

    \item \textbf{Mapeos SIFTS incompletos}: Entradas del archivo SIFTS donde
    \texttt{PDB\_BEG} es nulo (55\% de las entradas SIFTS) requieren inferir la
    correspondencia de numeraci\'on, lo que puede introducir errores de mapeo
    similares a los descritos anteriormente.
\end{itemize}

\subsection{Sesgo muestral y diversificaci\'on}

El dataset validado con SIFTS presenta un sesgo significativo: solo 91 prote\'inas
\'unicas representan los 5.115 pares validados inicialmente. Prote\'inas como la
lisozima (P00698, 777 pares) y la $\beta$-2-microglobulina (P61769, 672 pares)
dominan el dataset.

Para abordar este sesgo, se diversific\'o el dataset seleccionando un PDB
representativo por prote\'ina. De las 10.000 prote\'inas objetivo, se obtuvieron
10.432 con procesamiento exitoso. La mediana de RMSD subi\'o de \SI{0,64}{\angstrom}
a \SI{1,26}{\angstrom}, indicando que la precisi\'on del dataset validado estaba
parcialmente inflada por la sobrerepresentaci\'on de prote\'inas bien predichas.

\subsection{Implicaciones para la investigaci\'on}

\begin{enumerate}
    \item \textbf{Los modelos AlphaFold2 son confiables como punto de partida}:
    Con RMSD mediano de \SI{0,64}{\angstrom}, los modelos son suficientemente
    precisos para la mayor\'ia de aplicaciones, incluyendo docking molecular,
    dise\~no de mutantes y an\'alisis de interacciones.

    \item \textbf{El mapeo correcto es esencial}: Cualquier comparaci\'on
    estructural debe utilizar SIFTS u otro sistema de mapeo de posiciones.

    \item \textbf{Precauci\'on con conformaciones}: Los modelos AlphaFold2
    representan la conformaci\'on apo. Para la conformaci\'on holo (con ligando),
    verificar contra estructuras experimentales relevantes.

    \item \textbf{Validaci\'on por prote\'ina}: El RMSD promedio global
    enmascara variabilidad entre prote\'inas individuales.
\end{enumerate}

\subsection{Comparaci\'on con la literatura}

Nuestros resultados con mapeo correcto son consistentes con la literatura:

\begin{itemize}
    \item \citet{jumper2021highly} reportaron GDT-TS mediano $> 90$ en CASP14,
    equivalente a RMSD $< \SI{1}{\angstrom}$. Nuestro RMSD mediano de
    \SI{0,64}{\angstrom} es consistente.

    \item \citet{akdel2022structural} encontraron que el 36\% de los modelos
    ten\'ian RMSD $< \SI{1}{\angstrom}$. Nuestro 70,2\% es superior, posiblemente
    porque nuestro dataset incluye estructuras disponibles durante el entrenamiento.

    \item \citet{mariani2013lddt} establecieron que lDDT $> 0,7$ indica modelos
    de calidad \'util. Los modelos AlphaFold2 con RMSD $< \SI{2}{\angstrom}$
    superan ampliamente este umbral.
\end{itemize}

\subsection{Limitaciones del estudio}

\begin{enumerate}
    \item \textbf{Sesgo temporal}: Las prote\'inas comparadas probablemente
    estaban disponibles durante el entrenamiento de AlphaFold2.

    \item \textbf{Solo prote\'inas con estructura experimental}: Excluye
    prote\'inas dif\'iciles de cristalizar.

    \item \textbf{Comparaci\'on r\'igida}: El RMSD tras superposici\'on
    r\'igida no captura la calidad para prote\'inas con dominios m\'oviles.

    \item \textbf{SIFTS como \'unica fuente}: Puede contener errores o
    mapeos incompletos para algunas entradas PDB.
\end{enumerate}
